\section*{Step 47 Door-Test Statement}

Let \(q_{\rm QM}\) and \(q_{\rm GR}\) be the two access quotients of the
finite toy:
\[
q_{\rm QM}=(d_0,d_1,d_2), \qquad q_{\rm GR}=(d_0,d_2,d_3).
\]
The candidate joint quotient \(L=(d_0,d_1,d_2,d_3)\) has zero descent
residuals to both children.

The door-test asks whether a framework law forces the existence of such a
shared carrier. The answer on the tested laws is negative. F37 classifies
joint access versus complementarity; F51 gives the form of unification as a
common refinement; FoEC gives each theory a carrier but not one shared carrier;
and SAU licenses a non-descending object once its descents are audited. None
forces the existence of the common carrier for the QM/GR pair.

Verdict: the common-carrier premise is LANDED * GROUND (conditional), grounded
as the down-shadow of the semiclassical co-sourcing layer. This is a
recognition-mode landing, not an unconditional derivation and not merely
relocation.

The recognition source is semiclassical co-sourcing: one field \(\psi\) carries
both the Born readout \(|\psi|^2\) and the stress-energy readout \(T[\psi]\).
Step26 extends this warrant toward matter-sourced geometry by computing a
back-reaction loop \(V[\psi]=V_0+\kappa T_{00}[\psi]\). The non-co-sourced
control fails this warrant. Thus the common-carrier premise is warranted, not
arbitrary, but it is not an SBT-internal theorem. The remaining residual is the
free gravitational sector and trans-semiclassical frame transfer.
